\section{Existence theorem}

Let us first state \EGA III 3.4.2 in a slightly more
general form.

Let $f:\formelX \to \formelX'$ be an adic morphism\sfootnote{This hypothesis is not essential, \cf \Ref{XII}, p\ptbl\pageref{XII.14}.} of formal preschemes, with $\formelX'$ noetherian. Let $\Ii'$ be an ideal of definition of $\formelX'$; since $f$ is adic, $f^*\Ii'=\Jj$ is\nde{by the very definition, \cf \EGA I~10.12.1.} an ideal of definition of $\formelX$.

For every $n\in \NN$, set \label{IX.2}
\begin{equation} \label{eq:IX.2.1}
\formelX_n=(\formelX, \Oo_\formelX/{\Jb}^{n+1}),
\end{equation}
this is an ordinary prescheme which has the same underlying topological space as $\formelX$.

Let $\formelF$ be a \emph{coherent} $\Oo_\formelX$-Module. For every $k\in\NN$, the $\Oo_{\formelX_k}$-Modules
\begin{equation} \label{eq:IX.2.2}
F_k=\formelF/{\Jb}^{k+1}\formelF
\end{equation}
are \emph{coherent}. For every $i$, one has a homomorphism
\begin{equation} \label{eq:IX.2.3} \psi_i: R^if_*(\formelF) \to \varprojlim_k R^if_*(F_k),
\end{equation}
deduced by functoriality from the natural homomorphism:
\begin{equation} \label{eq:IX.2.4} \formelF \to F_k.
\end{equation}

Set
\begin{equation} \label{eq:IX.2.5}
\bS= \gr_{{\Ib}}\Oo_{\formelX'}=\tbigoplus_{k \in \NN} {\Ib}^k/{\Ib}^{k+1},
\end{equation}
\begin{equation} \label{eq:IX.2.6} \gr_{\Jb}(\formelF)=\tbigoplus_{k \in \NN} {\Jb}^k\formelF/{\Jb}^{k+1}\formelF,
\end{equation}
\begin{equation} \label{eq:IX.2.7} \KK^i=R^if_*(\gr_{{\Jb}}(\formelF))=\tbigoplus_{k \in \NN}R^if_* ({\Jb}^k \formelF/{\Jb}^{k+1}\formelF\;).
\end{equation}
It is clear that $\KK^i$ is equipped with a structure of graded $\bS$-Module.

\begin{theorem} \label{IX.2.1}
Suppose that $\KK^i$ is a graded $\bS$-Module \emph{of finite type} for $i=n-1$, $i=n$, $i=n+1$, then:
\begin{enumeratei}
\item
$R^nf_*(\formelF)$ is coherent.
\item
The homomorphism $\psi_n$ (\Ref{eq:IX.2.3}) is a topological isomorphism. The natural filtration of the second member of (\Ref{eq:IX.2.3}) is ${\Ib}$-good.
\item
The projective system of the $R^nf_*(F_k)$ satisfies the uniform Mittag-Leffler condition.
\end{enumeratei}
\end{theorem}

The proof is very easy from $\EGA 0_{\textup{III}} 13.7.7$ (\cf \EGA III $3.4.2$), provided one corrects the text on page 78 as indicated\refstepcounter{toto}\label{pagerectif} in (\EGA III~2, $\rm{Err}_{\rm{III}}24$).

\setcounter{equation}{7}
\refstepcounter{subsection}\label{IX.2.2}
\begin{enonce*}{Theorem 2.2\ndemark}
Let $A$ be a noetherian adic ring and let $I$ be an ideal of definition of $A$. Let $T$ be a closed subset of $X'=\Spec(A)$. Suppose that $I$ is generated by some $t \in A$. Let us resume the notations \Ref{eq:IX.1.31}, \Ref{eq:IX.1.32} and \Ref{eq:IX.1.33}. Let $\formelF$ be a coherent $\Oo_{\hat{X}}$-Module. Set \ndetext{many algebraization statements have been obtained since, not to mention those cited below, \cf the articles of Faltings or of Mme Raynaud cited in the editor's note~\eqref{HL}~p\ptbl\pageref{HL} and \eqref{MiR}~p\ptbl\pageref{MiR} respectively. One thinks in particular of the results of Artin (see in particular Artin~M., ``Algebraization of formal moduli. {I}'', in \emph{Global Analysis (Papers in Honor of K\ptbl Kodaira)}, {Univ. Tokyo Press}, {Tokyo}, {1969}, p\ptbl 21--71), but also of the recent results on algebraicity of leaves of foliations; see in particular Bost~J.-B., ``Algebraic leaves of algebraic foliations over number fields'', \emph{Publ. Math. Inst. Hautes \'Etudes Sci.} \textbf{93} (2001), p\ptbl 161--221 and Chambert-Loir~A. ``Algebraicity theorems in Diophantine geometry (after J.-B. Bost, Y. Andr\'e, D. \& G. Chudnovsky)'', in \emph{S\'eminaire Bourbaki, Vol. 2000/2001},  Ast\'erisque~, vol.~282, Soci\'et\'e math\'ematique de France, Paris, 2002, \Exp 886, p\ptbl 175--209 and the references cited. In particular, one will find in these two articles discussions of the link between algebraization questions and the theory of Diophantine approximation.}
\begin{equation} \label{eq:IX.2.8}
F_0=\formelF / {\Jb} \formelF,
\end{equation}
where ${\Jb}=t \Oo_{\hat{X}}$ is an ideal of definition of $\hat{X}$. \emph{Suppose that} $A$ is a quotient of a regular noetherian ring and that:
\begin{enumerate}
\item
$t$ is $\formelF$-regular,
\item
$\prof_{T'}F_0\geq 2$.
\end{enumerate}
\noindent
Then there exists a coherent $\OX$-Module $F$ such that $\hat{F}\simeq \formelF$.
\end{enonce*}

It suffices to prove that $\hat{f}_*(\formelF)$ is a coherent $\Oo_{\hat{X}}$-Module, where $\hat{f}: \hat{X} \to \hat{X'}$ is the morphism of formal preschemes deduced from the injection of $X$ into $X'$ by completion with respect to $t$. Indeed, $A$ is separated and complete for the $t$-adic topology, there will therefore exist an $A$-module $F'$ whose completion will be isomorphic to $\hat{f}_*(\formelF)$. Since $X$ is an open subset of $X'$, one may take $F=\tilde{F'}_{|X}$.

It remains to show that \Ref{IX.2.1} is applicable to the morphism of formal preschemes $\hat{f}$ and to $\formelF$. Now, by hypothesis (1), for every $k\in \NN$ one has an isomorphism:
\begin{equation*} {\Jb}^k \formelF /{\Jb}^{k+1}\formelF \to \formelF / {\Jb} \formelF,
\end{equation*}
whence it follows that the hypothesis of \Ref{IX.2.1} will be satisfied if one knows that
\[
R^if_*(F_0)\text{ is coherent for }i \leq 1.
\]
Now this follows from (2) and from the finiteness theorem \Ref{VIII.2.1}. Whence the conclusion.

It remains to specialize this statement by assuming that $A$ is a local ring.

\begin{corollary}\setcounter{equation}{8} \label{IX.2.3}
Let $A$ be a noetherian local ring and let $t \in \rr(A)$ be an element of the radical of $A$. Suppose that $A$ is separated and complete for the $t$-adic topology, and, moreover, a quotient of a regular ring (for example, suppose that $A$ is a complete noetherian local ring). Set
\begin{equation} \label{eq:IX.2.9} X'=\Spec A, \quad T = \{ \rr(A) \},
\end{equation}
and let us resume the notations (\Ref{eq:IX.1.31}), (\Ref{eq:IX.1.32}) and (\Ref{eq:IX.1.33}). Let $\formelF$ be an $\Oo_{\hat{X}}$-Module. Suppose that:
\begin{enumerate}
\item
$t$ is $\formelF$-regular,
\item
$\prof_{T'} F_0 \geq 2$, with $F_0=\formelF/{\Jb}\formelF$ and ${\Jb}=t\Oo_{\hat{X}}$.

Then there exists a coherent $\OX$-Module $F$ such that $\hat{F}\simeq \formelF$.
\end{enumerate}
\end{corollary}

Let us remark that here $T'$ is the set of prime ideals $\pp$ of $A$ such that \hbox{$\dim A/\pp =1$}.
